\chapter{Controle Externo}

\section{Controle da Administração Pública}

Controle é a atividade de verificar se a atuação estatal observa normas, objetivos, resultados, limites de competência e padrões de responsabilidade. Em um Estado Democrático de Direito, a Administração não atua sem supervisão: decisões, receitas, despesas, patrimônio, políticas e resultados podem ser submetidos a mecanismos de controle institucional e social.

O controle pode incidir antes, durante ou depois da atuação administrativa. Pode examinar conformidade jurídica, desempenho, integridade, qualidade da gestão, confiabilidade de informações e responsabilização.

\subsection{Classificações do controle}

\begin{table}[ht]
\centering
\small
\begin{tabularx}{\textwidth}{l X}
\toprule
Critério & Classificações usuais \\
\midrule
Momento & prévio/preventivo, concomitante e posterior. \\
Órgão controlador & administrativo, legislativo, judicial, controle externo e social. \\
Posição institucional & interno ou externo à estrutura controlada. \\
Objeto & legalidade, legitimidade, economicidade, desempenho, patrimônio, finanças, operações. \\
\bottomrule
\end{tabularx}
\end{table}

A mesma atividade pode receber mais de uma classificação. Uma auditoria realizada por Tribunal de Contas durante execução de contrato é controle externo e concomitante.

\section{Controle interno}

Controle interno é exercido dentro da própria estrutura estatal. A Constituição exige sistemas de controle interno em cada Poder e atribui-lhes funções de avaliação, acompanhamento e apoio ao controle externo.

Controle interno não deve ser confundido com uma única unidade administrativa. O \textbf{sistema} envolve políticas, responsabilidades, procedimentos e estruturas; uma controladoria ou unidade de auditoria interna pode integrar esse sistema, mas não o esgota.

\subsection{Funções do controle interno}

Entre suas funções estão:
\begin{itemize}
\item acompanhar cumprimento de metas e programas;
\item avaliar legalidade e resultados da gestão;
\item controlar operações de crédito, avais e garantias, conforme o regime aplicável;
\item apoiar o controle externo;
\item comunicar irregularidades relevantes aos órgãos competentes.
\end{itemize}

\section{Controle externo}

Controle externo é exercido por instituição ou Poder distinto da estrutura diretamente responsável pela gestão controlada, nos termos constitucionais.

No plano federal, o Congresso Nacional exerce controle externo com auxílio do Tribunal de Contas da União. Nos Estados, o modelo constitucional é reproduzido com Assembleia Legislativa e Tribunal de Contas estadual, observadas as peculiaridades constitucionais e legais.

\begin{teoria}[auxílio não é subordinação]
O Tribunal de Contas auxilia o Poder Legislativo, mas possui competências próprias diretamente previstas na Constituição. Por isso, não é correto tratá-lo como mero órgão técnico subordinado hierarquicamente ao Parlamento.
\end{teoria}

\section{Controle social e accountability}

Controle social é exercido pela sociedade por meio de transparência, participação, ouvidorias, conselhos, denúncias, representações e acesso a informações públicas.

\termo{Accountability} descreve uma relação em que agentes públicos devem explicar e justificar sua atuação e podem sofrer consequências quando descumprem deveres. Ela possui, portanto, dimensão informacional e dimensão de responsabilização.

Prestação de contas é uma manifestação de accountability, mas não se resume à entrega de documentos. Uma prestação de contas adequada deve permitir compreender o uso dos recursos, a responsabilidade dos agentes e os resultados alcançados.

\section{Objeto constitucional da fiscalização}

O art. 70 da Constituição estabelece fiscalização contábil, financeira, orçamentária, operacional e patrimonial quanto à legalidade, legitimidade, economicidade, aplicação das subvenções e renúncia de receitas.

\subsection{Fiscalização contábil}

Examina registros, demonstrações, critérios e informações contábeis. Busca verificar se representam adequadamente os fatos e se observam normas aplicáveis.

\subsection{Fiscalização financeira}

Analisa fluxos financeiros, arrecadação, pagamentos, disponibilidades, endividamento e demais operações que envolvem recursos públicos.

\subsection{Fiscalização orçamentária}

Verifica execução de receitas e despesas em relação ao orçamento aprovado e às regras fiscais e orçamentárias.

\subsection{Fiscalização operacional}

Examina processos, governança e resultados. Pode avaliar economicidade, eficiência, eficácia e efetividade.

\subsection{Fiscalização patrimonial}

Abrange bens, direitos e obrigações, incluindo aquisição, guarda, uso, alienação e conservação do patrimônio público.

\section{Legalidade, legitimidade e economicidade}

Esses três critérios não são sinônimos.

\termo{Legalidade} pergunta se a conduta está de acordo com o ordenamento. \termo{Legitimidade} amplia a análise para finalidade pública, razoabilidade e aderência aos fins constitucionais. \termo{Economicidade} examina se recursos foram obtidos e utilizados em condições economicamente adequadas.

\begin{exemplo}[legal mas antieconômico]
Uma compra pode cumprir formalmente todas as etapas legais e, ainda assim, apresentar preço excessivo ou especificação desnecessariamente onerosa. O controle externo pode examinar economicidade sem precisar afirmar que todo problema econômico constitui ilegalidade formal.
\end{exemplo}

\section{Eficiência, eficácia e efetividade}

\begin{table}[ht]
\centering
\small
\begin{tabularx}{\textwidth}{l X}
\toprule
Conceito & Pergunta central \\
\midrule
Economicidade & os insumos foram obtidos em condições adequadas de custo e qualidade? \\
Eficiência & os recursos foram transformados em produtos com boa relação entre entradas e saídas? \\
Eficácia & as metas e produtos previstos foram alcançados? \\
Efetividade & os resultados produziram impacto real sobre o problema público? \\
\bottomrule
\end{tabularx}
\end{table}

Uma política pode ser eficiente na entrega de serviços e pouco efetiva se esses serviços não alterarem o problema social que justificou sua criação.

\section{Dever de prestar contas}

A Constituição determina que prestará contas qualquer pessoa física ou jurídica, pública ou privada, que utilize, arrecade, guarde, gerencie ou administre dinheiros, bens e valores públicos ou pelos quais o poder público responda, nos termos constitucionais.

O dever decorre da relação com recursos públicos, não apenas da condição formal de servidor.

\section{Tribunais de Contas: natureza e posição institucional}

Tribunais de Contas são órgãos constitucionais de controle externo. Não integram o Poder Judiciário e não exercem jurisdição judicial, embora pratiquem julgamentos administrativos de contas e adotem decisões com efeitos jurídicos relevantes.

Suas competências derivam diretamente da Constituição e são complementadas por leis orgânicas e regimentos.

\section{Competências constitucionais dos Tribunais de Contas}

No paradigma federal, o art. 71 da Constituição atribui ao TCU conjunto amplo de competências.

\subsection{Parecer prévio sobre contas do chefe do Executivo}

O Tribunal \textbf{aprecia} as contas anuais do chefe do Poder Executivo e emite parecer prévio. Esse parecer subsidia julgamento pelo Poder Legislativo competente.

\subsection{Julgamento de contas de administradores e responsáveis}

O Tribunal \textbf{julga} contas dos administradores e demais responsáveis por dinheiros, bens e valores públicos e de quem der causa a perda, extravio ou irregularidade de que resulte prejuízo ao erário.

\begin{atencao}[apreciar e julgar]
A distinção verbal é essencial: contas de governo do chefe do Executivo recebem parecer prévio; contas de responsáveis por gestão de recursos são julgadas pelo Tribunal, respeitado o regime constitucional aplicável.
\end{atencao}

\subsection{Atos sujeitos a registro}

O Tribunal aprecia, para fins de registro, legalidade de atos de admissão de pessoal e concessões de aposentadorias, reformas e pensões nas hipóteses constitucionais, observadas exceções e jurisprudência.

\subsection{Auditorias e inspeções}

Pode realizar auditorias e inspeções por iniciativa própria ou por provocação dos legitimados constitucionais. Auditoria é trabalho sistemático baseado em critérios, evidências e procedimentos. Inspeção tende a ter escopo mais delimitado e objetivo específico.

\subsection{Sanções}

Tribunais de Contas podem aplicar sanções previstas em lei. O exercício sancionador exige tipicidade, competência, processo regular, contraditório, motivação e individualização de responsabilidade.

\subsection{Determinações}

Quando identifica ilegalidade, o Tribunal pode fixar prazo para que órgão ou entidade adote providências necessárias ao cumprimento da lei, nos termos constitucionais e legais.

\subsection{Sustação}

A Constituição distingue sustação de ato e tratamento de contrato. A sustação de ato impugnado pode ser realizada pelo Tribunal nos termos constitucionais. Em contratos, há disciplina própria que envolve inicialmente o Congresso Nacional no modelo federal, com atuação subsequente do Tribunal em determinadas condições.

\section{Contas de governo e contas de gestão}

\termo{Contas de governo} avaliam macrogestão do chefe do Executivo: planejamento, execução orçamentária, resultados fiscais e balanço geral. \termo{Contas de gestão} examinam atos de administração de recursos por responsáveis sujeitos ao julgamento técnico.

A distinção não deve ser aplicada mecanicamente sem considerar jurisprudência e desenho institucional do ente, mas é base importante para compreender parecer prévio e julgamento de contas.

\section{Prestação de contas, tomada de contas e tomada de contas especial}

\termo{Prestação de contas} é apresentada pelo responsável de forma ordinária, nos termos definidos. \termo{Tomada de contas} pode ser instaurada para suprir ausência ou apurar situação quando a prestação regular não ocorre adequadamente.

\subsection{Tomada de contas especial}

Tomada de contas especial é processo de natureza excepcional destinado, em linhas gerais, a apurar fatos, identificar responsáveis e quantificar dano quando houver omissão no dever de prestar contas ou indícios de dano ao erário e os mecanismos administrativos ordinários não solucionarem a situação.

Ela não deve ser banalizada como primeiro instrumento de apuração. Em geral, busca-se antes adotar medidas administrativas para recompor dano e esclarecer fatos.

\section{Responsabilização no controle externo}

A responsabilização exige vincular pessoa, conduta e consequência jurídica.

\subsection{Conduta e nexo causal}

Não basta que um dano tenha ocorrido durante a gestão de determinado agente. É preciso demonstrar relação entre sua ação ou omissão e o resultado imputado, considerando suas competências e circunstâncias.

\subsection{Dano ao erário}

Dano deve ser demonstrado e, quando cabível, quantificado. Irregularidade formal não implica automaticamente prejuízo financeiro.

\begin{exemplo}[irregularidade sem dano]
Um contrato pode apresentar falha formal em documento de fiscalização sem que isso, por si só, prove pagamento indevido. A irregularidade pode justificar determinação ou sanção, conforme o caso, mas imputação de débito exige demonstração do prejuízo e de responsabilidade correspondente.
\end{exemplo}

\subsection{Responsabilidade solidária}

Solidariedade não é presumida. Deve decorrer de fundamento jurídico e da demonstração de que os agentes se enquadram nos pressupostos que justificam responsabilidade conjunta.

\section{Processo de controle externo}

Decisões de Tribunais de Contas resultam de processo estruturado. Embora detalhes variem entre tribunais, aparecem fases como autuação, instrução, manifestações técnicas, contraditório, parecer do Ministério Público de Contas quando cabível, julgamento e recursos.

\subsection{Contraditório e ampla defesa}

Quando decisão pode impor débito, sanção ou outra consequência desfavorável, o responsável deve ter oportunidade de conhecer imputações e apresentar defesa, salvo situações cautelares excepcionais disciplinadas pelo ordenamento.

\subsection{Citação e audiência}

Em muitos regimes, \termo{citação} é utilizada quando há débito ou imputação mais grave, enquanto \termo{audiência} solicita razões sobre irregularidade sem débito. A terminologia exata depende da lei orgânica e do regimento do Tribunal competente.

\subsection{Motivação da decisão}

A decisão precisa indicar fundamentos jurídicos e elementos probatórios relevantes. Relatório técnico é parte da instrução, mas não substitui automaticamente a deliberação do órgão competente.

\section{Evidência de controle}

Evidência é o conjunto de informações utilizadas para sustentar achados e conclusões. Deve ser suficiente em quantidade e apropriada em qualidade e pertinência.

Documentos oficiais, registros de sistemas, bases de dados, inspeções, confirmações externas, entrevistas e análises podem compor a evidência. Uma única declaração verbal, sem corroboração, geralmente possui força probatória menor que evidências independentes e verificáveis.

\section{Achado de auditoria}

Uma estrutura didática comum para achado inclui:
\begin{itemize}
\item \textbf{critério}: o que deveria ocorrer;
\item \textbf{condição}: o que foi encontrado;
\item \textbf{causa}: por que ocorreu;
\item \textbf{efeito}: consequência ou risco;
\item \textbf{evidência}: suporte factual da constatação.
\end{itemize}

Nem todo trabalho exige apresentar esses componentes com os mesmos rótulos, mas a lógica ajuda a produzir conclusões verificáveis.

\section{Determinações e recomendações}

\termo{Determinação} ordena providência vinculada a dever jurídico e deve possuir fundamento claro e objeto executável. \termo{Recomendação} propõe melhoria quando existe espaço legítimo para escolha gerencial.

Uma recomendação excessivamente prescritiva pode invadir competência do gestor. Uma determinação vaga pode ser impossível de monitorar.

\section{Medidas cautelares}

Medida cautelar busca preservar utilidade da decisão final e evitar dano, agravamento ou ineficácia do processo. Possui natureza instrumental, não sancionadora.

Requisitos variam conforme o regime, mas normalmente envolvem plausibilidade jurídica e risco decorrente da demora.

\begin{atencao}
Cautelar não antecipa punição. Seu objetivo é proteger o processo ou o patrimônio até decisão definitiva.
\end{atencao}

\section{Controle preventivo e concomitante}

Controle preventivo ocorre antes da consumação do ato; controle concomitante acompanha execução em andamento. Ambos podem evitar dano, mas há um limite institucional importante: o órgão de controle não deve substituir o gestor nem assumir responsabilidade pela decisão administrativa.

O Tribunal pode apontar ilegalidades e riscos, mas não deve se tornar coautor das escolhas discricionárias legítimas da gestão.

\section{Monitoramento}

Monitoramento verifica implementação de determinações e recomendações e, quando pertinente, seus efeitos.

Cumprimento formal não significa necessariamente resolução do problema. Se a recomendação buscava reduzir indisponibilidade de sistema, instalar ferramenta de monitoramento pode ser uma ação implementada; ainda é necessário verificar se o controle funciona e se o risco foi efetivamente reduzido.

\section{Denúncia e representação}

Denúncias e representações são formas de levar fatos ao conhecimento do Tribunal. A existência de denúncia não comprova a irregularidade. Ela fornece informação inicial que deve ser submetida a critérios de admissibilidade e instrução.

\section{Prescrição e segurança jurídica}

A atividade sancionadora e de ressarcimento deve observar regimes constitucionais, legais e jurisprudenciais de prescrição. Fórmulas antigas de imprescritibilidade genérica não devem ser usadas sem verificar a natureza da pretensão e a jurisprudência vigente.

Segurança jurídica também impõe atenção a estabilidade de decisões, proteção da confiança e efeitos do tempo sobre relações administrativas.

\section{Controle externo de Tecnologia da Informação}

Em TI, o objeto de controle pode ser muito mais amplo do que conferir aquisição de equipamentos. Exemplos:
\begin{itemize}
\item governança e alinhamento estratégico;
\item gestão de riscos e segurança;
\item contratações de software, nuvem e serviços;
\item continuidade e recuperação;
\item gestão de identidades e acessos;
\item qualidade e integridade de dados;
\item sistemas que processam receitas e despesas;
\item algoritmos e inteligência artificial;
\item proteção de dados pessoais;
\item níveis de serviço e desempenho;
\item resultados de transformação digital.
\end{itemize}

\begin{exemplo}[auditoria de sistema de arrecadação]
O auditor não deve limitar-se a verificar se o contrato do sistema foi licitado. Pode examinar controles de acesso, segregação de funções, trilhas de auditoria, integridade dos cálculos, gestão de mudanças, disponibilidade, conciliação com registros contábeis e tratamento de incidentes.
\end{exemplo}

\section{Síntese conceitual}

\mapafuncional{Controle Externo}
{Controle externo = critérios + evidência + responsabilização + melhoria da gestão pública}
{Base constitucional $\rightarrow$ fiscalização contábil, financeira, orçamentária, operacional e patrimonial.\\Competências $\rightarrow$ parecer prévio, julgamento, registro, auditoria, sanção e determinação.\\Processo $\rightarrow$ instrução, contraditório, decisão e recursos.\\Auditoria $\rightarrow$ critério, condição, causa, efeito e evidência.\\Pós-decisão $\rightarrow$ monitoramento.}
{Chefe do Executivo: parecer prévio.\\Gestor de recursos: julgamento técnico.\\Débito: demonstre dano + nexo + responsabilidade.\\Cautelar: preserve utilidade, não puna.\\Recomendação: preserve espaço legítimo de gestão.}
{Auxílio ao Legislativo $\neq$ subordinação.\\Ilegalidade $\neq$ dano automático.\\Denúncia $\neq$ prova.\\Cautelar $\neq$ sanção.\\Controle concomitante $\neq$ cogestão.}
{Qual é o objeto?\\Quem é o responsável?\\Qual é o critério?\\Que evidência sustenta a condição?\\Há dano quantificado?\\A medida é corretiva, sancionadora ou cautelar?}

\begin{resumo}
Controle externo deve ser compreendido como sistema de accountability constitucional. A questão central não é apenas encontrar irregularidades, mas relacionar competência, critério, evidência, responsabilidade e consequência jurídica. Para Auditor/TI, essa lógica será aplicada a sistemas, dados, contratos, segurança, continuidade e governança tecnológica.
\end{resumo}
